\documentclass[11pt,a4paper]{article}
\usepackage[margin=1in]{geometry}
\usepackage{booktabs}
\usepackage{longtable}
\usepackage{graphicx}
\usepackage{hyperref}
\usepackage{listings}
\usepackage{xcolor}
\usepackage{enumitem}
\usepackage{caption}
\usepackage{amsmath}
\usepackage{amssymb}
\usepackage{tcolorbox}
\tcbuselibrary{breakable}

\hypersetup{
    colorlinks=true,
    linkcolor=blue,
    filecolor=magenta,
    urlcolor=cyan,
}

\tcbset{
    promptbox/.style={
        colback=blue!5!white,
        colframe=blue!75!black,
        title={#1},
        fonttitle=\bfseries,
        before skip=6pt,
        after skip=6pt,
        breakable,
    }
}

% Long unbreakable \texttt{} paths are the main source of overfull boxes in
% this report; a generous emergencystretch plus sloppy paragraph breaking lets
% TeX loosen those lines instead of running into the margin.
\setlength{\emergencystretch}{3em}
\sloppy

\title{Performance Report 29: Vulkan Engine \\
\large Before the First Frame --- Setup Cost and the Octree Load Path}
\author{}
\date{\today}

\begin{document}

\maketitle

\begin{abstract}
Reports~20--28 optimized what the engine draws. This report turns to what
happens before the first interactive frame. In the recorded run
\texttt{LocalScene::loadScene} took \textbf{81.31\,s}
(\texttt{logs/run.log:315}) and the first \texttt{drawFrame()} reported
\textbf{81.92\,s} (\texttt{logs/run.log:325}). The chunk-topology pass
accumulated \textbf{83.89\,s} of worker time over 1024 chunks
(\texttt{logs/run.log:762}), but its own cumulative averages carry the tell:
39.3\,ms/chunk for the first 256, 164.6 for the second, 119.6 for the third,
and 4.2 for the last 256 (\texttt{run.log:371,660,750,762}). The per-chunk
work did not get forty times cheaper; the three layer workers spent the middle
of the load queued behind each other. Callgrind attributes \textbf{237.8M}
corner-cache probes to \textbf{32.1M} real SDF evaluations (86\% hits),
15.5M \texttt{shape} calls, 5.45M node allocations and about 3.8M task
enqueues. The audit finds \textbf{9 issues} (2 critical, 2 high, 3 medium,
2 low): a float-keyed corner cache whose leaf pass runs twice, a per-chunk
two-submit topology path serialized on a global mutex, two debug walks that
run with their flags off and retain roughly 478\,MB, and allocation churn in
the hot shape loop. Scope score \textbf{4.3/10}; folded into the combined
estimate at 10\% weight, \textbf{5.8/10}.
\end{abstract}

\tableofcontents
\newpage

% =====================================================================
\section{Executive Summary}
% =====================================================================

\subsection{Scope and Method}

Scope is bring-up from process start to the first non-stalling frame. That
covers the setup entry (\texttt{MyApp::setup}, \texttt{MyApp.cpp:749}), scene construction
(\texttt{MainSceneLoader::loadScene}, \texttt{utils/MainSceneLoader.hpp:310}),
\texttt{LocalScene::loadScene} (\texttt{utils/LocalScene.cpp:563}), the 315
octree \texttt{apply} calls they issue, the chunk-topology flatten/compute
path (\texttt{services/OctreeTopologyCompute.cpp}), the two debug collection
walks in \texttt{SceneRenderer::processNodeLayer}
(\texttt{vulkan/renderer/SceneRenderer.cpp:2205}), and the upload/submit
plumbing they share. Out of scope, and deliberately so: the rendering
optimizations of reports~20--28, GDAL heightmap load, texture-array decode,
shader compilation and pipeline creation, input, and the ImGui layer.

Method matches the earlier reports: full reads with file:line evidence.
The application was \emph{not} re-executed for this report. Wall-clock numbers
come from the recorded native run \texttt{logs/run.log} (clean tree at
\texttt{f6c5687b}); instruction and call counts come from
\texttt{logs/callgrind.out}, whose \texttt{loadScene} took 3912.75\,s
(\texttt{logs/callgrind.log:1801}) --- a 48$\times$ inflation that serializes
the three layer workers. Instruction shares are therefore used only for
\emph{relative} attribution, never as wall-clock. Struct sizes used in the
memory arithmetic were verified by compiling the real headers:
\texttt{CubeSDF} = 80\,B, \texttt{CubeWithColor} = 64\,B,
\texttt{NodeOperationResult} = 120\,B, \texttt{OctreeThreadContext} = 240\,B.
\texttt{VULKAN\_SETUP\_DIAG} was not enabled in the recorded run (no
\texttt{[setup-diag]} lines), which is itself issue M7.

\subsection{The Gate: Mesh-Identical}
\label{sec:gate29}

Same three qualifying classes as the rendering reports, translated to setup:

\begin{enumerate}[leftmargin=1.5em,itemsep=2pt]
\item The \emph{same pure SDF evaluation} --- same function, same point, same
float arithmetic --- executed fewer times. Independent pure evaluations may be
reordered or moved between threads; no surviving evaluation may change its
inputs or its result.
\item \emph{Provably dead work} removed: a walk whose output is never read, a
store never loaded, a pass whose values are already in hand.
\item \emph{Identical bytes through a cheaper path}: the same node tables, the
same dispatches in the same per-chunk order, batched into fewer submissions or
one upload arena.
\end{enumerate}

Every prompt fails closed on a mesh-identical A/B: byte-identical
\texttt{OctreeNode} SDF arrays and \texttt{OctreeTopologyInput}; identical
triangle set after canonical sort (the packed pool's vertex order is not
reproducible even run-to-run by design,
\texttt{shaders/computation/octree/OctreeTriangles.comp:799--801}); and
pixel-identical overlays for the debug paths. No change to \texttt{minSize},
chunk size, the LOD ladder, thread counts, or scene content is admitted by
this gate.

% The budget table is too tall to share a page with the gate list above it;
% starting it on a fresh page keeps the heading, table and caption together
% instead of leaving the heading orphaned above a deferred float.
\clearpage
\subsection{Measured Budgets}

\begin{table}[h]
\centering
\small
\begin{tabular}{@{}p{6.1cm}rp{5.0cm}@{}}
\toprule
\textbf{Measurement} & \textbf{Value} & \textbf{Evidence} \\
\midrule
\texttt{LocalScene::loadScene} wall time & 81.31\,s & \texttt{run.log:315} \\
First \texttt{drawFrame()} wall time & 81.92\,s & \texttt{run.log:325} \\
Topology worker time, 1024 chunks & 83.89\,s cumulative & 1024 $\times$ 81.926\,ms, \texttt{run.log:762} \\
Topology rate, chunks 1--256 & 39.3\,ms/chunk (10.07\,s) & \texttt{run.log:371} \\
Topology rate, chunks 257--512 & 164.6\,ms/chunk (42.14\,s) & \texttt{run.log:660} \\
Topology rate, chunks 513--768 & 119.6\,ms/chunk (30.62\,s) & \texttt{run.log:750} \\
Topology rate, chunks 769--1024 & 4.2\,ms/chunk (1.07\,s) & \texttt{run.log:762} \\
Long frames while chunks streamed & 115, 126, 148\,ms & \texttt{run.log:476,479,738} \\
SDF corner-cache probes & 237.77M & callgrind calls \\
Real SDF distance evaluations & 32.1M (86\% hits) & callgrind calls \\
\texttt{Octree::shape} invocations & 15.52M & callgrind calls \\
\texttt{buildShapeSDF} invocations & 29.72M & callgrind calls \\
Node allocations & 5.45M & callgrind calls \\
Task enqueues (debug / shape) & 2.27M / 1.47M & callgrind calls \\
Node-table upload traffic & $\sim$289\,MB & 2.256M nodes $\times$ 128\,B \\
Debug caches retained, flags off & $\sim$478\,MB & 5.15M$\times$80\,B + 1.03M$\times$64\,B \\
Texture arrays, main & 720 / 747\,MB & \texttt{run.log:323} \\
\bottomrule
\end{tabular}
\caption{All wall-clock values are from the native run; all counts are from
callgrind. Cumulative worker time exceeds the wall because three layer
workers overlap. The 4.2\,ms/chunk of the last window is the uncontended
cost; the first 768 chunks average 107.8\,ms, a $\sim$26$\times$ convoy
factor. If the 81.92\,s first frame is the topology convoy --- M7 exists to
prove it --- the three workers achieved a parallel efficiency of only 1.02$\times$
(83.89\,s of work in 81.92\,s).}
\end{table}

\subsection{Banked Before Report 29}
\label{sec:banked29}

The setup path is not starting from zero, and the report does not re-propose
what already works. Banked: the GPU-resident topology path itself --- only
node SDF tables cross to the GPU, \texttt{OctreeTriangles.comp} writes the
packed pools directly, and only integer counters return
(\texttt{OctreeTopologyCompute.cpp:360--367}); device-local keys and dedup
table (\texttt{:169--186}); the span-relative pool binding that stays inside
every device's \texttt{maxStorageBufferRange} (\texttt{:269--299}); the
\texttt{hostReadBarrier} discipline (\texttt{:196--208}); and the first
\texttt{VULKAN\_SETUP\_DIAG} hooks (\texttt{MyApp.cpp:754--772}). Report~28's
items are all rendering and do not overlap this scope.

\subsection{Recommended Order}

H3 (gate the dead debug walks) $\rightarrow$ C2 (batch the topology
submissions) $\rightarrow$ C1 (one leaf pass, integer corner keys)
$\rightarrow$ H4 (task and map churn) $\rightarrow$ M7 (per-phase timers)
$\rightarrow$ M5 (pool topology) $\rightarrow$ M6 (descriptor writes)
$\rightarrow$ L9 (key headroom) $\rightarrow$ L8 (batched wakeups).
H3 is two conditions and removes 60\% of the task traffic; C2 is the largest
wall-clock item; C1 is the largest instruction item. M7 is listed before
M5/M6 because it is what makes their deltas measurable, not because it is
large.

% =====================================================================
\section{Cost Model}
% =====================================================================

\begin{equation}
\begin{aligned}
C_{\mathrm{setup}} = &\underbrace{A\,(C_{\mathrm{apply}} + N_s\,C_{\mathrm{shape}} + E\,C_{\mathrm{probe}})}_{\text{octree construction}} \\
&+ \underbrace{\sum_{c=1}^{F} \bigl( S\,(C_{\mathrm{submit}} + C_{\mathrm{fence}}) + U_c\,C_{\mathrm{upload}} + K\,C_{\mathrm{desc}} \bigr)}_{\text{chunk topology}} \\
&+ \underbrace{D\,(C_{\mathrm{sdfwalk}} + C_{\mathrm{bboxwalk}})}_{\text{debug collection}} + \underbrace{T\,C_{\mathrm{task}} + M\,C_{\mathrm{alloc}}}_{\text{allocation churn}}
\end{aligned}
\end{equation}

Measured values: $A = 315$ applies (35 direct sites plus 24 road segments
and 256 ring segments, \texttt{MainSceneLoader.hpp:498,522});
$N_s = 15.52$M \texttt{shape} calls; $E = 237.8$M probes of which 32.1M miss;
$F = 1024$ chunks; $S = 2$ submissions per chunk (2048 fence waits);
$U_c = 128$\,B/node, 289\,MB total; $K = 26$ descriptor writes per chunk;
$D = 2$ walks per chunk; $T \approx 3.8$M tasks; $M = 5.45$M node
allocations. C1 removes the second leaf pass and cuts $C_{\mathrm{probe}}$;
C2 attacks $S$, $U_c$ and $K$; H3 removes $D$ when the overlays are off; H4
attacks $T$ and the per-node maps. The equation's outputs --- the node SDF
arrays and the triangle set --- are unchanged for every item: the gate.

% =====================================================================
\section{Memory Residency}
% =====================================================================

The scene's texture arrays are the dominant resident allocation and are not
this report's target: 720 of 747\,MB, five maps at 149\,MB each
(\texttt{run.log:323}). The octree itself is 5.45M nodes at 72\,B
($\sim$392\,MB) plus allocator block overhead; that is the scene, not churn,
and H4 leaves it alone.

The avoidable resident cost is the debug caches: 5.15M \texttt{CubeSDF}
records at 80\,B (412\,MB) plus 1.03M \texttt{CubeWithColor} records at 64\,B
(66\,MB) --- 478\,MB retained for overlays whose flags default to
\texttt{false} (\texttt{types/Settings.hpp:101--102}). H3 removes them when
the overlays are off and frees them when the toggle goes off.

The topology buffers are small in residency but large in traffic: the seven
host-visible input buffers are reused and sized to the largest chunk (a few
MB), while \texttt{uploadInput} writes 128\,B per node per chunk --- 289\,MB
across the load. C2 keeps the same bytes but moves them through one staging
arena per layer instead of 1024 mapped writes.

Transient churn is larger than either: 124.18M \texttt{NodeOperationResult}
constructions at 120\,B is 14.9\,GB of constructor traffic (stack-scoped, not
resident), 26M \texttt{std::unordered\_map} constructions in the child loop,
and 3.8M futures with their shared state. H4 addresses all three without
changing a single SDF value.

% =====================================================================
\section{Critical Issues}
% =====================================================================

\subsection{C1: The Leaf Corner Pass Runs Twice, and the Cache Is Float-Keyed}
\label{sec:c1}

\textbf{Severity: Critical.}

Two exact problems share one loop. First, every leaf evaluates its eight
corners twice. \texttt{shape} calls
\texttt{buildShapeSDF(\ldots,true)} unconditionally at
\texttt{space/Octree.cpp:456}; for a leaf, that fills \texttt{r.shapeSDF}
from the eight corners at \texttt{:272--275}. At \texttt{:597} the function
calls \texttt{buildShapeSDF} again, and because \texttt{isLeaf} is true the
same branch re-evaluates the same eight points with the same
\texttt{cubeMin}/\texttt{cubeLength} (\texttt{:272}). The cache turns the
second pass into hits, but a hit is still a probe: the counts are 29.72M
\texttt{buildShapeSDF} calls against 15.52M \texttt{shape} calls, so
$\sim$14.2M second calls produce up to $\sim$114M redundant probes --- 48\%
of the 237.8M total. Deleting the leaf arm of \texttt{:597} is exact: the
values are already in \texttt{r.shapeSDF}, and \texttt{buildResultSDF}
(\texttt{:290--307}) only reads them.

Second, the surviving probes are expensive. \texttt{evaluateSDF}
(\texttt{:257--265}) keys a \texttt{tsl::robin\_map} by a 12-byte
\texttt{glm::vec3}, hashed through \texttt{hashCombine(murmurMix(\ldots))}
and compared with a three-float equality. The profile's single largest
self-cost function is \texttt{glm::vec3(float,float,float)} at 78.75e9
instructions (6.78\% of the program), and the probe path adds
\texttt{robin\_hash::find\_impl} (22.17e9, 1.91\%), \texttt{operator==} on
\texttt{vec3} (9.79e9), \texttt{bucket\_entry::value} (9.32e9), and the hash
machinery (roughly 7e9 each). Together the probe path is on the order of a
tenth of all retired instructions in the unoptimized debug build.

The key can be exact and cheap because every corner of the recursion's own
leaves lies on the tree's lattice. \texttt{expand} doubles the root until it
contains the shape (\texttt{Octree.cpp:222--255}), so node lengths are
$L_{\mathrm{root}}/2^l$ and a leaf corner is \texttt{rootMin + n * L\_leaf},
where $L_{\mathrm{leaf}}$ is the leaf cube length (the first
$L_{\mathrm{root}}/2^l$ not exceeding \texttt{minSize}); leaf centers are
half steps of the same lattice. A packed \texttt{uint64} lattice index
(21 bits per axis) selects exactly the same entries as the float key whenever
the rounded index reconstructs the queried float bit-for-bit, which
\texttt{evaluateSDF} checks per query. A flat open-addressed table of 16-byte
entries, reserved to the previous apply's peak, replaces the node-based robin
map and its per-entry allocations. Queries that do not reconstruct --
accumulated subdivision rounding, or cubes finer than $L_{\mathrm{leaf}}$ that
the recursion follows because the stored tree holds nodes from an earlier
apply with a smaller \texttt{minSize} -- use the float-keyed fallback, so the
fast path is never a correctness requirement. (An early debug assert that the
division was integral was removed after it aborted at runtime: the
mixed-\texttt{minSize} applies in \texttt{MainSceneLoader} create exactly
those finer nodes, and they are a legitimate state, not an error.)

\begin{tcolorbox}[promptbox={Prompt --- C1: one leaf pass, integer corner keys}]
Goal: halve the probes, then make each probe cheap. Steps: (1) at
\texttt{Octree.cpp:597}, call \texttt{buildShapeSDF} only when
\texttt{!isLeaf}; add a debug assertion that \texttt{r.shapeSDF} equals a
freshly evaluated set for leaves; (2) add \texttt{rootMin} and the leaf
spacing $L_{\mathrm{leaf}}$ to \texttt{ShapeArgs}, key the cache by a packed
\texttt{uint64} lattice index, and accept a rounded index only when it
reconstructs the queried float bit-for-bit (\texttt{rootMin} + n*spacing);
anything else -- off-lattice cubes from stored finer nodes, or subdivision
rounding -- falls back to the float key; (3) replace
\texttt{robin\_map<vec3,float>} in \texttt{OctreeThreadContext} with a flat
open-addressed table reserved per apply; (4) keep the per-task contexts (they
capture the 86\% reuse) and add hit/miss counters under M7. Files:
\texttt{space/Octree.cpp}, \texttt{space/Octree.hpp},
\texttt{space/OctreeThreadContext.hpp}, \texttt{space/ShapeArgs.hpp}.
Validation: byte-identical \texttt{OctreeNode} SDF arrays for the fixed
scene, identical triangle set, probe count down by $\sim$48\%,
validation-clean.
\end{tcolorbox}

\subsection{C2: Chunk Topology Pays Two Global Fence Waits per Chunk}
\label{sec:c2}

\textbf{Severity: Critical.}

\texttt{computeResident} takes a process-wide mutex for its entire body
(\texttt{OctreeTopologyCompute.cpp:370}) and then calls
\texttt{runSingleTimeCommands} twice (\texttt{:409} and \texttt{:484}). That
helper takes a second global mutex for the whole record-submit-wait sequence
(\texttt{vulkan/core/VulkanApp.cpp:1060}) and waits on the fence while
holding it (\texttt{:1148}). For 1024 chunks that is 2048 global lock
acquisitions and 2048 fence waits, all on the graphics queue, with three
layer workers that cannot overlap. If the first-frame stall is this convoy
(M7 will confirm), the measured parallel efficiency is 1.02$\times$: 83.89\,s
of cumulative worker time against an 81.92\,s frame. The per-window rates
make the convoy visible --- 39.3, 164.6, 119.6,
then 4.2\,ms/chunk --- while the average node count falls monotonically
(3410, 3076, 2474, 2203), so the middle window is not doing more work. The
last 256 chunks are the queue with nothing in front of it.

Each chunk also uploads 128\,B per node through seven host-visible coherent
buffers (\texttt{:154--161}, \texttt{:225--250}): 2.256M nodes, 289\,MB, plus
a counter reset per chunk. It rewrites all 13 descriptor bindings twice
(\texttt{:403}, \texttt{:475}), 26.6k updates, although only bindings 11/12
change between the calls. None of this is wrong; it is simply unbatched.

The exact fix is per-layer batching. Flatten every chunk of a layer, append
its node table to one staging arena, and record all of that layer's walk
dispatches into one command buffer (one preallocated descriptor set per
chunk, binding per-chunk slices of the arena). Submit once; read all counters
with one map. Reserve each chunk's span from the exact counts, then record
all fill/dedup/index sequences into a second command buffer and submit once.
The result is 6 submissions and 6 fence waits for the whole scene (3 layers
$\times$ 2 phases) instead of 2048, three uploads instead of 1024, and 13k
descriptor writes instead of 26.6k. The per-chunk dispatch order, barriers
and shader code are unchanged; only the grouping and the arena offsets move.
The measured floor is 4.2\,ms/chunk, so 1024 chunks over three workers is
$\sim$1.4\,s of worker time even if batching saves nothing else.

\begin{tcolorbox}[promptbox={Prompt --- C2: batched topology, one arena}]
Goal: one submit per layer per phase, one upload per layer. Steps: (1) add a
per-layer arena holding all chunks' node tables and a counter array of four
uints per chunk; add a \texttt{nodeBase} push-constant field and rebase
\texttt{cellIds} by it (integer addition only); (2) flatten all chunks, append,
record walk dispatches with per-chunk descriptor sets into one command
buffer, submit once; (3) map the counters once, reserve spans from the exact
counts, record fill/dedup/index per chunk into a second buffer, submit once;
(4) keep the old per-chunk path behind a debug flag for A/B. Files:
\texttt{services/OctreeTopologyCompute.cpp/.hpp},
\texttt{services/OctreeTopologyTypes.hpp},
\texttt{shaders/\allowbreak computation/\allowbreak octree/\allowbreak OctreeTriangles.comp},
\texttt{vulkan/\allowbreak renderer/\allowbreak SceneRenderer.cpp}. Validation: mesh-identical
triangle set, synchronization validation plus
\texttt{VULKAN\_GPU\_ASSISTED=1 make run-debug}, validation-clean; check the
per-chunk fill/dedup barriers survive the merge.
\end{tcolorbox}

% =====================================================================
\section{High-Severity Issues}
% =====================================================================

\subsection{H3: Two Debug Walks Run With Their Flags Off}
\label{sec:h3}

\textbf{Severity: High.}

\texttt{SceneRenderer} always constructs both debug renderers
(\texttt{SceneRenderer.cpp:458--459}) and always initializes them
(\texttt{:563--564}). \texttt{processNodeLayer} therefore enters the SDF
cube block (\texttt{:2256--2279}) and the bounding-box block
(\texttt{:2287--2307}) for every chunk, and each block runs a full
\texttt{iterateMultiThreaded} over the chunk subtree
(\texttt{LocalScene.cpp:441,487}). The settings that are supposed to control
them default to \texttt{false} (\texttt{Settings.hpp:101--102}) and gate only
registration and drawing (\texttt{MyApp.cpp:1996--2001,2891,3348}), never
collection.

The cost is measured: 8.4M node visits and 5.15M cubes for the SDF walk,
13.5M visits and 1.03M boxes for the bbox walk; 2.27M of the 3.8M total task
enqueues (60\%) come from these two walks; the retained records are 478\,MB.
All of it is invisible in the default configuration. When the user does
enable an overlay, the walk is identical, so the gate holds.

The fix is an enabled bit on each renderer, set when the settings toggle
changes, plus a one-shot full rebuild when it turns on and a clear when it
turns off. The collection blocks then read the bit instead of the pointer.
Toggling on mid-session must repopulate every live chunk, not just newly
built ones; the existing per-chunk map makes that a dirty-mark pass.

\begin{tcolorbox}[promptbox={Prompt --- H3: gate debug collection on the settings}]
Goal: no walk, no retention, no tasks when the overlay is off. Steps:
(1) add \texttt{setEnabled(bool)} and an atomic flag to
\texttt{DebugSDFRenderer} and \texttt{DebugCubeRenderer}; (2) gate the two
blocks in \texttt{processNodeLayer} on the flag; (3) on toggle-on, clear and
mark all live chunks for a one-shot recollection; on toggle-off, clear the
stored records; (4) keep the draw-side settings checks as they are. Files:
\texttt{vulkan/\allowbreak renderer/\allowbreak SceneRenderer.cpp},
\texttt{vulkan/\allowbreak renderer/\allowbreak debug/\allowbreak DebugSDFRenderer.*},
\texttt{vulkan/\allowbreak renderer/\allowbreak debug/\allowbreak DebugCubeRenderer.*}, \texttt{MyApp.cpp}.
Validation: with both flags off, zero cubes/boxes and zero debug tasks
(counters); with a flag on, pixel-identical overlay A/B against the current
build; validation-clean.
\end{tcolorbox}

\subsection{H4: Per-Node Allocation Churn in the Shape Loop}
\label{sec:h4}

\textbf{Severity: High.}

The shape loop allocates in four places that do not need to. (1)
\texttt{NodeOperationResult} is 120\,B and is constructed 124.18M times: the
\texttt{children[8]} array at \texttt{Octree.cpp:449--455} (960\,B per
\texttt{shape} call, each element default-initialized to \texttt{INFINITY}),
the \texttt{childResult[8]} parameter, and temporaries. That is 14.9\,GB of
constructor traffic. (2) \texttt{std::unordered\_map<int,int> brushCounts}
is declared inside the eight-child loop at \texttt{:679}; with at most eight
children and at most eight distinct brush indices, it is a linear scan
wearing a hash table: one construction and destruction per child per
non-leaf node, roughly 15M in the profile. (3) \texttt{shapeChildren} builds a
\texttt{std::vector} of futures per call (\texttt{:357}) without a reserve;
use a fixed array of eight plus a count. (4) \texttt{ThreadPool::enqueue}
allocates a \texttt{make\_shared} packaged task and its future state for
every task (\texttt{space/ThreadPool.tpp:11--17}) --- 3.8M heap allocations.
In \texttt{shapeChildren} the futures are only waited on
(\texttt{Octree.cpp:419}); a detached enqueue plus the existing
\texttt{inFlightShapeOps} counter removes the future, the shared state and
the allocation. \texttt{IteratorHandler} genuinely needs the future for
work-stealing (\texttt{space/IteratorHandler.cpp:54}), so leave that path
alone --- and H3 removes most of its traffic anyway.

None of this touches arithmetic. The octree node allocator's 5.45M
allocations are the tree itself and stay.

\begin{tcolorbox}[promptbox={Prompt --- H4: remove per-node allocation churn}]
Goal: no heap allocation in the shape loop's child bookkeeping. Steps:
(1) replace \texttt{brushCounts} with an eight-entry stack array and a linear
scan, preserving the strict-greater tie rule; (2) replace the futures vector
with a fixed array of eight and a count, or with detached enqueues plus the
existing atomic counter; (3) leave \texttt{IteratorHandler}'s cooperative
futures untouched; (4) add a debug-only allocation counter to prove the drop.
Files: \texttt{space/Octree.cpp}, \texttt{space/ThreadPool.tpp/.hpp},
\texttt{space/IteratorHandler.cpp}. Validation: byte-identical node SDF
arrays and triangle set; validation-clean.
\end{tcolorbox}

% =====================================================================
\section{Medium-Severity Issues}
% =====================================================================

\subsection{M5: One-Level Fork Band, Three Pools, Fresh Contexts}

\textbf{Severity: Medium.}

\texttt{isThreadNode} admits exactly one ladder level:
\texttt{minSize*8 < L <= minSize*16} (\texttt{Octree.cpp:343--345}). A fork
therefore exposes at most eight tasks, and every threaded child constructs a
fresh \texttt{OctreeThreadContext} (\texttt{:408}) --- a 240-byte object plus
an empty hash map --- discarding the parent's cache at the fork. Three layer
workers, an \texttt{Octree} pool sized to
\texttt{hardware\_concurrency()} (\texttt{Octree.hpp:41}), a
\texttt{LocalScene} pool of the same size (\texttt{LocalScene.cpp:174}) and a
per-layer generation pool all exist at once; during setup most of those
threads are parked on the two global mutexes of C2.

Fork two levels (up to 64 tasks) to keep a shared pool fed, size that pool
once, and recycle contexts from a small per-thread pool so the corner cache
survives between tasks. The pure evaluations are unchanged; only which
thread computes a given corner changes.

\subsection{M6: Descriptor Rewrites Twice per Chunk}

\textbf{Severity: Medium.}

\texttt{writeDescriptors} fills all 13 bindings
(\texttt{OctreeTopologyCompute.cpp:301--337}) and is called twice per chunk:
once with a null target before the walk (\texttt{:403}) and once with the
span before dedup (\texttt{:475}). Only bindings 11/12 differ; the other
nine are rewritten 1024 times each for nothing. Either write the input
bindings once per capacity change and only 11/12 per chunk, or push the 13
descriptors per dispatch with \texttt{VK\_KHR\_push\_descriptor} when
supported. With C2's per-chunk sets, this becomes one write per chunk.

\subsection{M7: Setup Has No Per-Phase Instrumentation}

\textbf{Severity: Medium (blocking for verification, not for correctness).}

\texttt{VULKAN\_SETUP\_DIAG} wraps only the top-level setup steps
(\texttt{MyApp.cpp:754--772}) and was off in the recorded run.
\texttt{LocalScene::loadScene} times only its total
(\texttt{LocalScene.cpp:565--578}). \texttt{countStats} prints a cumulative
average every 256 chunks (\texttt{OctreeTopologyCompute.cpp:340--357}) ---
which is exactly what hid the 26$\times$ convoy: the cumulative average at
1024 is 81.9\,ms, while the last window is 4.2\,ms. There are no counters for
cache hits/misses, tasks, uploaded bytes, fence time, descriptor time, or
debug-walk time.

Add a small setup profiler with scoped timers and relaxed atomics, enabled by
the existing environment variable: per apply, flatten, upload, walk, fence,
descriptor and debug walk; cache hits/misses per apply; bytes and tasks per
chunk. Print a final table with the last-window rate as well as the average.
Instrumentation only; no behavior change.

% =====================================================================
\section{Low-Severity / Hygiene Issues}
% =====================================================================

\subsection{L8: One Wakeup per Task}

\texttt{ThreadPool::enqueue} calls \texttt{condition.notify\_one()} for every
task (\texttt{ThreadPool.tpp:25}). With 3.8M tasks that is 3.8M wakeups,
most of them while workers are already running. Notify once after a fork's
eight enqueues, or only when the queue transitions from empty to non-empty.
The \texttt{SmallFunction} inline wrapper (16-byte shared pointer,
\texttt{:22--23}) is already the right shape and stays.

\subsection{L9: The Key-Capacity Guess Is 12 Vertices Short}

\texttt{kVerticesPerCellHint} is 24
(\texttt{OctreeTopologyCompute.cpp:31}), which underestimates dense chunks.
The log records eight retries; two of them are
\texttt{emitted=275508 cap=275496} and \texttt{emitted=13308 cap=13296}
(\texttt{run.log:341--343,412,464,554,633}) --- short by 12 to 24 vertices.
Each retry doubles the key space, re-uploads the full node table and re-runs
the walk: a complete extra pass. Add a small headroom
(\texttt{cells * 25 + 64}) and keep the retry as a safety net. Capacity only;
the emitted keys are unchanged.

% =====================================================================
\section{Deliberately Unchanged}
\label{sec:unchanged29}

\begin{itemize}[leftmargin=1.4em,itemsep=2pt]
\item \textbf{Texture decode and GDAL heightmap.} 720/747\,MB of texture
arrays and the TIF load are I/O- and decode-bound; changing formats or
resolutions is a content change, not an exact win.
\item \textbf{Shader compilation and pipeline creation.} One-time and cached;
not the target of this report.
\item \textbf{Scene content.} The 315 applies (35 direct sites, 24 road segments,
256 ring segments) define the scene; reducing them changes what is built.
\item \textbf{\texttt{minSize}, chunk size and the LOD ladder.} Quality knobs.
\item \textbf{The root-doubling lattice.} C1's exactness proof depends on it;
it is asserted, not changed.
\item \textbf{\texttt{apply}'s write lock.} Correctness; C1 and C2 reduce the
work under it, not the lock.
\item \textbf{The octree node allocator.} 5.45M allocations are the tree, not
churn.
\item \textbf{The packed pool's vertex order.} Not reproducible by design, so
the gate compares the triangle set, never the byte order
(\texttt{OctreeTriangles.comp:799--801}).
\item \textbf{\texttt{runSingleTimeCommands} semantics.} Kept for one-shot
work; C2 batches around it rather than replacing it.
\item \textbf{Thread count defaults.} M5 changes pool topology and fork
granularity, not the configured count.
\item \textbf{Debug overlay content.} H3 gates collection, not which cubes or
boxes are drawn.
\end{itemize}

% =====================================================================
\section{Implementation Roadmap and Expected Deltas}
% =====================================================================

\begin{table}[h]
\centering
\small
\begin{tabular}{@{}lp{5.6cm}p{6.0cm}@{}}
\toprule
\textbf{Item} & \textbf{Change} & \textbf{Est.\ effect} \\
\midrule
C1 & One leaf pass; integer flat corner cache & $-114$M probes ($-48\%$); probe path $\sim$3$\times$ cheaper \\
C2 & Per-layer batched topology, one arena & 2048 $\rightarrow$ 6 submits/waits; 289\,MB via 3 uploads; 26.6k $\rightarrow$ 13k writes \\
H3 & Gate debug walks on the settings & $-2.27$M tasks ($-60\%$), $-22$M visits, $-478$\,MB \\
H4 & Stack brush counts, fixed futures, detached shape tasks & $-15$M maps, $-3.8$M futures/allocs, $-14.9$\,GB ctor traffic \\
M5 & Shared pool, two-level fork, context recycling & fewer parked/contending threads; cache survives forks \\
M6 & Write bindings 9--12 once, or push descriptors & $-9.2$k redundant writes \\
M7 & Setup profiler with last-window rates & per-phase attribution; makes the rest verifiable \\
L9 & Key headroom & $-8$ full extra walk passes \\
L8 & Batched notify & $-3.3$M wakeups \\
\bottomrule
\end{tabular}
\end{table}

C1--C2 are the structural wins; H3 is the cheapest high-severity item and
should land first. The measured floor for C2 is the last-window rate
(4.2\,ms/chunk), so the topology convoy has roughly a 26$\times$ headroom
before any GPU-side change is needed. C1's instruction share is large but its
wall-clock share is not yet attributable without M7; the estimate is
deliberately a range, not a promise. Nothing here changes a rendered pixel or
a node value.

Explicitly out of scope (carried forward): the rendering items of reports
20--28, GDAL heightmap load (\texttt{math/HeightMapTif.cpp}), input/events
(\texttt{events/EventManager.cpp}), and the ImGui widget layer (release-only,
never profiled).

% =====================================================================
% Keep the final score section on one page and in source order: the rubric
% table is set as a non-float so it cannot be flushed above its heading.
\clearpage
\section{Overall Score}
\label{sec:score29}
% =====================================================================

\subsection{Rubric (Report-29 Scope)}

\begin{center}
\small
\begin{tabular}{@{}lcp{7.4cm}@{}}
\toprule
\textbf{Category} & \textbf{Wt.} & \textbf{Score and justification} \\
\midrule
Octree evaluation and cache & 35\% & \textbf{4.5} --- the GPU-resident
output path and the cache's 86\% hit rate are good; the leaf pass runs twice,
the key is a 12-byte float vector, and the probe path is a tenth of the
profile. \\
Chunk topology build/submit & 30\% & \textbf{4.0} --- device-local keys and
dedup, span-relative bindings and clean barriers are good; two global fence
waits per chunk and a 1.02$\times$ parallel efficiency are not. \\
Load-path memory and allocation & 20\% & \textbf{3.5} --- pools are
preallocated and the output never round-trips; 478\,MB of dead debug caches,
289\,MB of host-visible uploads and 26M maps drag it down. \\
Instrumentation and hygiene & 15\% & \textbf{5.5} --- the setup-diag hooks and
the topology stats exist; per-phase timing, cache counters and last-window
rates do not. \\
\bottomrule
\end{tabular}
\end{center}

\subsection{Overall: 4.3/10 Scope, 5.8/10 Combined}

\begin{equation}
0.35(4.5) + 0.30(4.0) + 0.20(3.5) + 0.15(5.5) = \mathbf{4.3}
\end{equation}

Report~28 scored the combined engine 6.0/10. Re-measuring this scope:

\begin{equation}
0.90(6.0) + 0.10(4.3) = \mathbf{5.8}
\end{equation}

The setup path has never had the equivalent of reports~20--28. It is not
broken --- the GPU-resident topology output is a genuinely good design ---
but it is unbatched, uninstrumented, and paying for two debug overlays that
are off. C1--C2--H3 are the items that would move it; the combined score
rounds down a fifth of a point because the recorded run spent roughly 163\,s
before a non-stalling frame, and none of that time is currently attributable
to a phase.

\end{document}
